\documentclass[11pt]{article}
\usepackage[margin=1in]{geometry}
\usepackage{amsmath,amssymb}
\usepackage{hyperref}
\usepackage{listings}
\usepackage{xcolor}
\usepackage{booktabs}
\usepackage{natbib}

\lstset{
  basicstyle=\ttfamily\small,
  breaklines=true,
  columns=fullflexible,
  frame=single,
  backgroundcolor=\color{gray!8}
}

\title{\texttt{unitTransformDist}: A Distribution-Agnostic Toolkit for\\
Unit-Interval Models via the Odds-Inverting Transform}

\author{
  Mohamad Yamen Al Mohamad \and
  Ali Akbar Heydari \and
  Hossein Jabbari Khamnei\\
  \small Department of Statistics, University of Tabriz, Tabriz, Iran
}

\date{}

\begin{document}
\maketitle

\begin{abstract}
\noindent
Data restricted to the open interval $(0,1)$ -- proportions, rates,
probabilities -- are commonly modelled in R with \texttt{betareg} (beta
regression), \texttt{simplexreg} (simplex regression), or the recently
released \texttt{gkwdist} (generalized Kumaraswamy family). Each of these
packages targets one specific parametric family. This paper introduces
\texttt{unitTransformDist}, which instead constructs a $(0,1)$-valued
distribution from an arbitrary base distribution on $(0,\infty)$ via the
transform $Y = 1/(1+X)$, taking the base distribution's density, CDF,
quantile, and random-generation functions directly as arguments rather
than selecting from a fixed registry. This design lets users plug in R's
built-in distributions or their own custom research distributions --
including ones with no closed-form quantile function -- without
modifying the package. Beyond this generality, we document and
empirically correct a subtle but consequential implementation error
common to profile-likelihood confidence intervals: freezing nuisance
parameters at their unconstrained maximum-likelihood estimate while
profiling a parameter of interest, rather than re-optimizing them at
each candidate value. A reproducible Monte Carlo study included in the
package ($n=60$, 150 replications, a two-parameter gamma base
distribution) shows this mistake reduces actual 95\% coverage to
approximately 59\%, and that the corrected procedure restores coverage
to approximately 94\%. The package also provides Monte Carlo utilities
for studying estimator bias and several bias-correction methods
(jackknife, bootstrap, Firth-type). We describe the package's design,
its relationship to existing CRAN alternatives, and its current
limitations.
\end{abstract}

\section{Introduction}

Continuous data confined to $(0,1)$ arise across many applied fields:
disease prevalence, vote shares, budget allocations, and normalized
indices are all commonly analyzed as proportions. The dominant approach
in R is beta regression via \texttt{betareg} \citep{cribarineto2010betareg},
which models the mean and precision of a beta-distributed response.
Alternatives include \texttt{simplexreg} for the simplex distribution
\citep{zhang2016simplexreg} and the recent \texttt{gkwdist} package,
which implements the seven-member generalized Kumaraswamy family with
C++-accelerated likelihood and derivatives
\citep{carrasco2010generalized}. A related but architecturally different
idea appears in \texttt{Newdistns}, which takes an arbitrary baseline
CDF $G$ and layers a fixed menu of generator families on top of it
(Marshall--Olkin $G$, exponentiated $G$, beta $G$, and related
constructions).

\texttt{unitTransformDist} addresses a narrower transform -- the
odds-inverting map $Y = 1/(1+X)$, the inverse of the classical odds
transform -- but removes the restriction to a fixed set of base
distributions entirely: the base distribution's d/p/q/r functions are
supplied directly by the caller. We do not claim this transform itself
is new; our contribution is (a) a general-purpose implementation that
works unmodified with any base distribution the user can write in R,
including one with no closed-form quantile function, and (b) a
documented and empirically verified fix to a profile-likelihood
implementation mistake that we believe is easy to make and likely
present, in some form, in other unit-interval fitting code.

\section{Design considerations}

The package targets applied statisticians who already have a candidate
base distribution in mind -- possibly a distribution not available in
any existing package -- and want to fit its unit-transform version
without first registering it anywhere. Three design decisions follow
from this.

First, \texttt{utd()} accepts \texttt{ddist}, \texttt{pdist},
\texttt{qdist}, and \texttt{rdist} as plain R functions, matching the
base-R d/p/q/r convention. \texttt{qdist} and \texttt{rdist} are
optional: if omitted, quantiles are obtained by numerically inverting
\texttt{pdist} with \texttt{stats::uniroot()}, and random draws are
obtained by inverse-CDF sampling through that same numerical quantile
function. This lets a user supply only a density and a CDF for an
entirely custom distribution and still get a fully functional object.

Second, every exported, user-facing function validates its inputs at
the boundary (domain checks, named-vector checks, \texttt{NA} checks)
and raises an informative error via \texttt{stop()}, rather than
propagating \texttt{NaN} silently. Internal helpers (prefixed with a
dot) skip this validation, since their callers have already validated
it.

Third, the package uses base R's S3 object system and implements the
generics a user already expects: \texttt{print()}, \texttt{summary()},
\texttt{coef()}, \texttt{logLik()}, \texttt{vcov()}, \texttt{AIC()},
\texttt{BIC()}. \texttt{summary()} methods return an object of a
\texttt{summary.*} class rather than printing directly, and a
corresponding \texttt{print.summary.*} method does the printing --
matching the convention used by \texttt{summary.lm()} in base R, and
avoiding \texttt{cat()}/\texttt{print()} calls inside computational
functions.

\section{Typical workflow}

\begin{lstlisting}[language=R]
library(unitTransformDist)
d <- utd(dgamma, pgamma, qgamma, rgamma, params = list(shape = 2.5))
y <- rutd(500, d)
fit <- fit_utd_mle(y, dgamma, start = c(shape = 1), lower = c(shape = 1e-6))
fit
profile_likelihood_ci(fit, "shape")
\end{lstlisting}

A fully custom, unregistered distribution follows the identical pattern:

\begin{lstlisting}[language=R]
dcustom <- function(x, rate, log = FALSE) dexp(x, rate = rate, log = log)
pcustom <- function(q, rate, lower.tail = TRUE) pexp(q, rate, lower.tail = lower.tail)
d2 <- utd(dcustom, pcustom, params = list(rate = 1.5))
qutd(0.5, d2)  # numerically inverted, no qdist supplied
\end{lstlisting}

\section{The profile-likelihood correction: technical details and
empirical evidence}
\label{sec:contribution}

For a model with parameter vector $\theta = (\theta_1,\dots,\theta_p)$,
the profile likelihood for $\theta_1$ at a candidate value $t$ is
defined as
\begin{equation}
  L_{\mathrm{prof}}(t) \;=\; \max_{\theta_2,\dots,\theta_p} \; L(t,\theta_2,\dots,\theta_p),
\end{equation}
i.e., the nuisance parameters must be re-optimized, not fixed, at every
candidate $t$. An implementation that instead evaluates the likelihood
at $(t,\hat\theta_2,\dots,\hat\theta_p)$ -- the nuisance parameters
frozen at their unconstrained MLE -- computes something else entirely,
and the resulting ``confidence interval'' from inverting a $\chi^2$
threshold against this quantity has no guaranteed coverage properties.

\texttt{unitTransformDist::profile\_likelihood\_ci()} performs the
correct re-optimization at every grid point via nested calls to
\texttt{stats::optim()}. We also provide
\texttt{safe\_profile\_likelihood()}, a version of the same routine
that is deliberately factored out to accept any user-supplied negative
log-likelihood function, independent of the rest of the package, so the
fix can be audited or reused on its own.

To quantify the practical impact,
\texttt{compare\_profile\_likelihood\_coverage()} runs both versions --
naive fixed-nuisance and corrected -- on the same simulated datasets and
reports empirical 95\% coverage:

\begin{lstlisting}[language=R]
true_dist <- utd(dgamma, pgamma, qgamma, rgamma,
                  params = list(shape = 2.5, rate = 1))
compare_profile_likelihood_coverage(
  true_dist, c(shape = 2.5, rate = 1), dgamma,
  start = c(shape = 1, rate = 1), param_name = "shape",
  n = 60, n_sim = 150, lower = c(shape = 1e-6, rate = 1e-6), seed = 1)
\end{lstlisting}

\begin{table}[h]
\centering
\begin{tabular}{@{}lcc@{}}
\toprule
Method & Coverage (95\% target) & $n_{\text{valid}}$ \\
\midrule
Naive (fixed nuisance) & 0.587 & 150 \\
Corrected (re-optimized) & 0.940 & 150 \\
\bottomrule
\end{tabular}
\caption{Empirical coverage of the naive versus corrected profile
likelihood, $n=60$, 150 Monte Carlo replications, gamma base
distribution with shape $=2.5$, rate $=1$.}
\label{tab:coverage}
\end{table}

The naive version's actual coverage (58.7\%) falls far short of its
nominal 95\% target; the corrected version (94.0\%) is close to
nominal. This is, to our knowledge, the first published empirical
quantification of this specific mistake's coverage impact for a
unit-interval model, and we present it as the package's central
contribution. We flag as experimental (and explicitly do not recommend
for production use) an additional analytic bias-correction routine,
\texttt{bias\_correct\_analytic()}, whose derivation exploits that the
transform's Jacobian is parameter-free, so the classical Cox and Snell
O(1/n) correction \citep{coxsnell1968} transfers unchanged from the base
distribution's likelihood; the current finite-difference implementation
of the required third derivatives has not yet been validated against a
closed-form case and should be treated as a documented but unverified
idea rather than a working tool. We report this limitation explicitly
rather than omitting it.

\section{Performance and limitations}

The package has no compiled code and depends only on \texttt{stats} and
\texttt{graphics}; every operation is pure R. This keeps the dependency
footprint minimal but means large-sample maximum likelihood fitting and
the Monte Carlo/bootstrap utilities scale linearly in the number of
\texttt{optim()} calls, which can be slow for very large
\texttt{n\_sim} or \texttt{n\_boot}; we have not benchmarked or
parallelized this. Where quantiles or random draws require numerical
inversion (\texttt{qdist}/\texttt{rdist} omitted), each call to
\texttt{qutd()}/\texttt{rutd()} performs one \texttt{uniroot()} search
per element, which is markedly slower than a closed-form quantile
function. The package currently provides no regression, mixture, or
censored-data extensions of the transform; these are natural next steps
but are not part of this release, and the current version's
mathematical verification effort has focused specifically on the
profile-likelihood correction rather than on every function in the
package.

\section{Discussion and conclusion}

\texttt{unitTransformDist} does not claim to be the first package to
generalize across base distributions (\texttt{Newdistns} already does
this for a different family of transforms), nor does it claim the
odds-inverting transform itself is new. Its value is narrower and, we
think, still useful: a minimal-dependency, any-distribution
implementation of this specific transform, paired with a concrete,
reproducible demonstration that a common profile-likelihood
implementation error meaningfully damages confidence interval coverage
on this model class. We see the distribution-agnostic
\texttt{safe\_profile\_likelihood()} routine as reusable beyond this
package's own scope, and we hope documenting this fix explicitly --
with the coverage numbers to show it matters -- is useful to other
package authors implementing profile likelihood for their own models.

\bibliographystyle{plainnat}
\bibliography{references}

\end{document}
